\clearpage % Start a new page
\lhead{\emph{Trabajos Relacionados}}
\chapter{Trabajos Relacionados}
\label{sec:Trabajos Relacionados}

En este capítulo, abordaremos los estudios y avances realizados a lo largo del tiempo en la predicción del nivel del río Paraguay. Primero, presentaremos las instituciones cuyo trabajo incluye el monitoreo y análisis del río, destacando sus proyectos y contribuciones en este ámbito. Luego, exploraremos enfoques estadísticos utilizados para predecir sus niveles. Finalmente, introduciremos un enfoque basado en Inteligencia Artificial, el cual constituye la base del modelo que presentamos en este libro.


\section{Instituciones y proyectos relacionados}

Es fundamental destacar las instituciones vinculadas a la hidrología en Paraguay, dado que el país depende en gran medida de sus ríos. A continuación, se presenta información sobre algunas de estas instituciones y sus roles en la gestión y estudio del río Paraguay.

\subsection{Ministerio del Ambiente y Desarrollo Sostenible (MADES)}

El \textit{Ministerio del Ambiente y Desarrollo Sostenible} (MADES) fue creado mediante el Artículo 3º de la Ley Nº 6123/18, que elevó al rango de ministerio a la antigua \textit{Secretaría del Ambiente}. Esta institución es la autoridad de aplicación de la Ley Nº 3239/07 \textit{“De los Recursos Hídricos del Paraguay”} \cite{MAD23}.

De acuerdo con la Ley Nº 3239/07, el MADES es responsable de la regulación y gestión sostenible e integral de todas las aguas y los territorios que las generan, sin importar su ubicación, estado físico o forma de ocurrencia dentro del territorio paraguayo. Su objetivo es garantizar el aprovechamiento social, económico y ambientalmente sostenible de los recursos hídricos para la población del país \cite{MAD23}.

Dentro del MADES, la gestión de los recursos hídricos está a cargo de la \textit{Dirección General de Protección y Conservación de los Recursos Hídricos} (DGPCRH). Según el Artículo 25 de la Ley Nº 1561/00, esta dirección tiene la función de formular, coordinar y evaluar políticas para la conservación de los recursos hídricos y sus cuencas, asegurando el mantenimiento de los caudales básicos, la capacidad de recarga de los acuíferos y el equilibrio entre los diferentes usos del agua \cite{MAD23}.

La DGPCRH cuenta con dos direcciones temáticas \cite{MAD23}:
\begin{enumerate}
    \item \textbf{Dirección de Hidrología e Hidrogeología}: Responsable del monitoreo hidrológico a nivel nacional (cantidad y calidad del agua), en coordinación con otras instituciones del Estado.
    \item \textbf{Dirección de Gestión de Cuencas Hidrográficas}: Encargada de coordinar, a nivel de cuencas hidrográficas, la participación de los usuarios, gobiernos nacionales y locales para la gestión sustentable del agua en el territorio nacional.
\end{enumerate}

\subsection{Centro Internacional de Hidroinformática de ITAIPU (CIH-ITAIPU)}

En julio del año 2006, ITAIPU Binacional y la \textit{United Nations Educational, Scientific and Cultural Organization} (UNESCO) firmaron un memorando de entendimiento, comprometiéndose a realizar esfuerzos conjuntos para el establecimiento del centro \cite{CIH23}.  

En junio de 2007, Brasil y Paraguay presentaron ante la UNESCO la propuesta del Centro Categoría 2, y en octubre del mismo año, el Director General de la UNESCO, en su informe 177 EX/66, recomendó la aprobación del CIH-ITAIPU como Centro Categoría 2 en la Asamblea General de la UNESCO, destacando el compromiso de ITAIPU para su implementación y consolidación \cite{CIH23}.  

En mayo de 2009, ITAIPU Binacional y la UNESCO firmaron la extensión del memorando de entendimiento, estableciendo directrices para la creación del CIH-ITAIPU con una validez hasta 2016 \cite{CIH23}.  

Finalmente, en marzo de 2016, en París, durante la 199.ª reunión del Consejo Ejecutivo de la UNESCO, se autorizó a la Directora General a firmar los acuerdos relativos a la creación del Centro Internacional de Hidroinformática para la Gestión Integrada de Recursos Hídricos en ITAIPU Binacional (Brasil y Paraguay), como un centro auspiciado por la UNESCO (Categoría 2) \cite{CIH23}.  

\subsubsection{Proyecto YRATO}

El acceso y la interpretación de datos hidrológicos en la sociedad representan un desafío para varios países donde las tecnologías de la información aún no tienen un protagonismo suficiente. Profesionales de distintas disciplinas, no solo de la comunidad hidrológica, requieren insumos relacionados con variables hidrométricas para sus estudios, enfrentando como principal inconveniente la dificultad de acceso a estos datos \cite{VV18}.  

La información se encuentra dispersa o simplemente no está disponible debido a la falta de tecnologías que faciliten su difusión. Además, las variables hidrológicas proporcionadas por entidades o estudios suelen presentarse de manera poco comprensible para el ciudadano no técnico, reduciendo su utilidad en la vida cotidiana \cite{VV18}.  

En respuesta a estas necesidades, el Centro Internacional de Hidroinformática desarrolló la aplicación \textit{YRATO}. En su primera etapa, esta herramienta permite monitorear en tiempo real variables hidrológicas como niveles hidrométricos, calidad del agua y precipitación satelital a nivel de la Cuenca del Plata, hasta la confluencia de los ríos Paraguay y Paraná \cite{yrato-link, VV18, CON20}.

\subsection{Dirección de Meteorología e Hidrología (DMH)}

Según el Decreto Ley N.º 9926 del 7 de noviembre de 1938, se creó la Dirección de Meteorología, reorganizando los Servicios Meteorológicos bajo una única Dirección, dependiente del Ministerio de Guerra y Marina, actualmente Ministerio de Defensa Nacional \cite{DMH23b}.  

Paraguay participó en la duodécima Conferencia de Directores de la Organización Meteorológica Internacional (OMI), que tuvo lugar en Washington en 1947. En esta conferencia se propuso la creación de la Organización Meteorológica Mundial (OMM), organismo especializado de las Naciones Unidas para la meteorología (tiempo y clima), que desde 1950 se convirtió en el portavoz autorizado sobre el estado y el comportamiento de la atmósfera terrestre, su interacción con los océanos, el clima que produce y la distribución resultante de los recursos hídricos. Paraguay fue miembro fundador de la OMM y ha estado representado permanentemente en ella por el Director de la DMH \cite{DMH23b}.  

En 1986, la Dirección de Meteorología pasó a denominarse, según la Ley N.º 1.228 del 29 de diciembre de 1986, como Servicio Nacional de Meteorología e Hidrología (SENAMHI), dependiente del Ministerio de Defensa Nacional \cite{DMH23b}.  

El 31 de marzo de 1990, mediante el Decreto 25/90, se creó la DINAC, y el SENAMHI, hasta entonces dependiente del Ministerio de Defensa Nacional, pasó a formar parte de esta entidad estatal. Se constituyó como una de sus Direcciones orientadas al servicio de la aviación nacional. La Ley N.º 73/90, que aprueba con modificaciones el Decreto N.º 25/90 del 26 de julio de 1990, estableció como principales cometidos de esta Dirección la administración de la red nacional de estaciones meteorológicas oficiales, la recopilación y procesamiento de datos provenientes de dichas estaciones, la promoción del estudio, desarrollo e investigación en Meteorología e Hidrología en todo el territorio nacional, en coordinación con las instituciones estatales competentes en la materia. Además, se definió como la única institución oficial autorizada a nivel nacional para la provisión de servicios y productos de información meteorológica e hidrológica \cite{DMH23b}.  

\subsubsection{Proyecto YSYRY}  

El proyecto \textit{YSYRY} \cite{ysyry-link}, identificado como PAR08801, fue desarrollado en el Departamento de Hidrología de la Gerencia de Climatología e Hidrología de la DMH-DINAC. Su ejecución estuvo a cargo de la Organización de Aviación Civil Internacional (OACI) \cite{ysyry-link}. 

Al momento de escribir este trabajo, se constató que el proyecto \textit{YSYRY} cuenta con modelos utilizados para la predicción del caudal del río, por un lado, y del nivel del río por otro.

\section{Trabajos Estadisticos}

Durante la década de 1980, se emplearon métodos de correlación lineal para estimar el comportamiento futuro del nivel del río Paraguay, basándose en un punto de control aguas arriba \cite{CON20}. Sin embargo, a nivel local, hasta la entrega del informe final detallado en \cite{CON20}, no se contaba con un sistema de pronóstico confiable, exceptuando las curvas de correlación entre estaciones hidrométricas del tramo paraguayo. Estas curvas proporcionaban un pronóstico de niveles aguas abajo con un alcance de uno a cinco días \cite{CON20}.

El sistema, desarrollado en la década de los ochenta por una misión técnica de las Naciones Unidas, fue utilizado por la Administración Nacional de Navegación y Puertos (ANNP) y la Dirección de Meteorología e Hidrología (DMH) de la Dirección Nacional de Aeronáutica Civil (DINAC) \cite{CON20}. De acuerdo con el “1er informe técnico de avances del proyecto Pronóstico de niveles y caudales del río Paraguay hasta 60, 90 días”, estos métodos solo permitían obtener una proyección lineal de la traslación de la onda de crecida. Por lo tanto, si se presentaban lluvias significativas en las cuencas intermedias, estas no eran consideradas en la ecuación de regresión lineal, lo que alteraba sustancialmente los resultados \cite{CON20}.

Un trabajo realizado por el Ing. J. L. Benza, de la Facultad de Ingeniería de la Universidad Nacional de Asunción, en 1998, propuso ajustes automáticos de distribuciones estadísticas para determinar las alturas hidrométricas (nivel del río) para diferentes periodos de recurrencia en el río Paraguay. Este estudio estuvo asociado al informe "Zonificación de áreas inundables del río Paraguay, CEN FIUNA (2000)" \cite{CON20}. Específicamente, el trabajo del Ing. Benza se centró en:

\begin{enumerate}
\item El ajuste de las posibles curvas de distribución a los datos observados de máximos anuales en varias estaciones.
\item La selección y utilización del modelo que mejor se ajuste a los datos observados en cada estación, para determinar las alturas en las respectivas estaciones.
\item El análisis de la posibilidad de estimar alturas faltantes (niveles de río faltantes) de algunas estaciones mediante la relación existente entre las alturas de estaciones vecinas.
\end{enumerate}

Una tesis de grado de 1999, realizada en la Facultad de Ingeniería de la Universidad Nacional de Asunción (FIUNA), actualizó la metodología del PNUD de 1980 \cite{CON20}.

Para la previsión de niveles en Asunción, la Dirección de Pronósticos Hidrológicos de la DMH-DINAC correlaciona puntos de esta misma estación considerando la influencia de los fenómenos “La Niña” y “El Niño” \cite{CON20}.

El informe detallado en \cite{CON20} indica que, para el pronóstico del nivel del río Paraguay, se analizó el comportamiento de correlaciones lineal, por un lado, y del modelo ARIMA por el otro. En el mencionado informe se explica que ARIMA mostró buenos resultados durante los meses en los que otros modelos presentaron dificultades.

Además, se especifica que el modelo ARIMA presentado es univariante, considerando como entrada únicamente los niveles observados y registrados en la estación de Asunción. Para meses como julio, agosto y septiembre, el modelo ARIMA presentó un horizonte de predicción incluso superior a un mes.

Una hipótesis mencionada en el informe respecto a la dificultad de previsión que presentan los modelos en el último trimestre del año es que podría estar relacionada con el efecto de la precipitación en estos meses, dado que representan el periodo de mayor precipitación anual.

\section{Trabajos con IA.}

En los últimos años, se ha observado un creciente interés en el uso de modelos basados en Inteligencia Artificial (IA) para la predicción del nivel del río Paraguay. Estos modelos, especialmente los basados en redes neuronales profundas (Deep Learning, DL), han demostrado ser herramientas valiosas para abordar las complejidades y no linealidades presentes en las series temporales hidrológicas.

Un estudio reciente aplicado a la predicción de niveles del río Paraguay en el puerto de Asunción comparó el rendimiento de modelos estadísticos y de DL utilizando datos de series temporales. Los resultados experimentales indicaron la superioridad de los modelos basados en DL, específicamente LSTM y GRU, sobre el modelo ARIMA, especialmente al incorporar niveles de otras estaciones como covariables y al realizar pronósticos con horizontes de 7, 14, 21 y 28 días \cite{amarilla2023comparative}. En comparación con los modelos estadísticos utilizados previamente, los modelos basados en DL redujeron significativamente el error cuadrático medio (Root Mean Squared Error, RMSE). Por ejemplo, para un horizonte de predicción de 28 días, el modelo GRU logró un RMSE de 60,9 cm, mejorando considerablemente respecto al modelo ARIMA, que obtuvo un RMSE de 77,1 cm. Esta mejora se observó consistentemente en todos los horizontes evaluados, demostrando la capacidad de los modelos de DL para capturar patrones más complejos y proporcionar predicciones más precisas.

Además, se destacó la importancia de incluir niveles de otras estaciones como covariables para mejorar la precisión y fiabilidad de los modelos de pronóstico. Esta inclusión permitió capturar patrones espaciotemporales y mejorar la robustez del modelo ante diferentes escenarios y condiciones cambiantes \cite{amarilla2023comparative}.

Por último, aunque los modelos de DL demostraron un mejor rendimiento, se sugiere que futuras investigaciones exploren la integración de estos modelos con enfoques basados en modelos físicos, lo que podría aprovechar el conocimiento específico del dominio y proporcionar predicciones aún más precisas y robustas \cite{amarilla2023comparative}.

